\begin{textsample}{NEG Dim 2 – am – Score -2.00 – am/am\_inf\_16.txt}  \label{ex:f2_neg_133}
Fotografia e Cinema Para além da possibilidade de se conversar com os bebês e crianças sobre as transformações de sua história a partir das fotografias pessoais, de família e de amigos, que trazem registros de passeios, festas e demais situações, as máquinas fotográficas ou \textbf{celulares} com \textbf{câmeras} possibilitam realizar imagens que colaboram com a nossa capacidade de olhar a partir de pontos de vistas alternativos : dependendo de como são realizadas as fotos, obtemos mais detalhes sobre os objetos focados, descobrindo particularidades que podem ainda não ter sido observadas.
Podemos nos colocar, junto às crianças, a observar variadas texturas, o que ocorre entre as formigas no parque da creche ou da pré-escola, os grãos de arroz no momento em que a refeição é servida, as paredes que delimitam as salas, as trilhas feitas na terra pelos insetos, os olhos, mãos e gestos dos amigos da turma, enfim, infinitas possibilidades de pesquisa nos ambientes vividos pelos meninos e meninas.
A educação do olhar é tão importante quanto a leitura das palavras.
A linguagem visual está nos diferentes espaços que a criança convive.
Mas que imagens as crianças têm disponível para exercitar modos de ver?
As fotografias de lugares e coisas não óbvias podem se constituir em repertório de imagens complexas e difíceis para desafiar o olhar das crianças, contribuindo para uma leitura mais ampla para além dos personagens de filmes e desenhos infantis que geralmente cobrem as paredes da escola da infância.
A princípio, é preciso que se tenham imagem de leitura fácil e imagens que instiguem o olhar a pensar sobre o que vê.
Esse exercício amplia a capacidade de imaginar e criar, além de garantir o uso de recursos tecnológicos e midiáticos pelas crianças, como asseguram as DCNEI ( BRASIL, 2009b ).
O cinema deve ser apresentado às crianças porque elas têm direito de acessar àquilo que de mais elaborado a humanidade já produziu.
É comum o uso de televisores nas escolas com vídeos de desenho animados, musicais e filmes de animação infantil, mas devemos oferecer, além disso.
Existe um grande acervo de curtas, documentários e filmes que não são animação que podem favorecer a ampliação de repertórios das crianças.
Importante destacar que o uso do vídeo não pode ser um passatempo sem planejamento usado para emergências ou deixar as crianças expostas por um longo tempo, mesmo quando não há interesse.
O vídeo deve ser visto antes pela/pelo professora/professor, sempre ser mediado por diálogos e ter objetivo de promover o desenvolvimento do olhar curioso e atento das crianças, ampliando seu conhecimento de mundo.

% matched lemmas: celular, câmera
\end{textsample}
